\subsection{Structural assumptions}
\label{sec:struct}

Table~\ref{tab:struct} lists the assumptions that are built into the simulator's rules rather than exposed as
parameters, with where each lives in the code and which results we judge most exposed to it. The last column is
our reasoning, not a test: none of these assumptions has been varied, so we do not know how large any of the
effects is.

{\small
\begin{longtable}{>{\raggedright\arraybackslash}p{1.45in}>{\raggedright\arraybackslash}p{2.6in}>{\raggedright\arraybackslash}p{1.9in}}
\caption{Structural assumptions embedded in the simulator. None is varied in any experiment.}
\label{tab:struct}\\
\toprule
Assumption & What the simulator does & Results most exposed\\
\midrule
\endfirsthead
\toprule
Assumption & What the simulator does & Results most exposed\\
\midrule
\endhead
\bottomrule
\endfoot
Students never leave the queue &
Every student who is not a designated walker waits for an elevator, however long the line. Nobody takes the stairs
on their own, gives up, or leaves early (\texttt{split\_walkers}). The walking measure is applied to a baseline in
which no one walks. &
Baseline wait and lateness; the size of the walking effect; the benefit of any measure that shortens lines.\\
Compliance is a per-car coin flip &
Backpack compliance $b$ adds exactly two seats to a car load with probability $b$, independently of the car, the
floor, and the riders (\texttt{fresh\_capacity}). Compliance is never clustered by class, floor, or time of day, and the
bonus does not depend on the base capacity or on how many people are waiting. &
The size of the backpack effect; its near-linearity in compliance; the ranking of backpacks over walking.\\
Space is the only constraint on boarding &
Everyone in line boards if there is room; nobody declines a crowded car or waits for the next one. Boarding is
first-in-first-out, and an empty car takes the direction of the first person in line (\texttt{serve\_floor}). &
Wait and the share of empty stops.\\
A stop costs only door time &
Travel is a constant 3~s per floor with no acceleration. Everything else about stopping is folded into a fixed
8~s plus 1~s per rider (\texttt{t\_door}, \texttt{t\_pax}). The door-time sweep varies the value of this cost, not its form. &
The cost of empty stops, hence the value of removing riders from lower floors; the gap between short and long trips.\\
Controller and queues &
Collective rule; one shared queue per floor from which any eligible car draws; the two-car bank serves lower floors on
the same terms as the four-car bank; no zoning, no priority for the upper floors. Cars start the period at random
floors, with no backlog from the previous period (\texttt{World}, \texttt{depart}). &
The finding that floor~10 is the worst floor; the offset results; absolute clearing times.\\
Demand is generated independently and in full &
Outgoing and incoming occupancy of a room are independent draws with the same $p_{occ}$. Every seat of every incoming
class is filled, and nobody is absent. Leavers depart uniformly in the first two minutes, with no early or late
dismissal; non-transferring arrivals come only through the lobby on a triangular schedule; transfers choose a room in
proportion to open seats (\texttt{build\_demand}). &
Queue size and shape; lateness; the effect of the offset.\\
Parameters tuned to the story &
$p_{occ}$ was chosen so that the baseline shows a long line without collapse, and the transfer share $q$ and
adjacent share $a$ are not measured. That the baseline resembles the description is therefore not evidence that the
model is right. &
All absolute levels.\\
Lateness is a hard threshold &
A student is late if they arrive after the nominal class start; there is no grace period and no instructor discretion.
Offset classes keep their length. &
Lateness percentages; the offset trade-off.\\
Every changeover is the same &
One 15-minute period, the same for the 50-minute and 75-minute patterns, and for every block of the day. &
Any statement about a particular time of day.\\
\end{longtable}
}

Two points follow. First, the sensitivity analyses in Section~\ref{sec:occ} vary parameters inside this fixed
structure. They show that the ordering of the measures is robust to occupancy, door time, and adjacent-floor
traffic. They do not show that it would survive different student behavior or a different controller. Second, the
effects are not all of one kind. Some assumptions, such as students never leaving the queue, make the baseline look
worse and could shrink the apparent benefit of every measure if they were relaxed. Others, such as the
coin-flip compliance, affect how much each measure achieves. We have not tested any of this, and we have not
tried to guess its net direction.
